\documentclass[10pt,a4paper]{article}

\usepackage{marktwo}
\usepackage{graphicx}

\begin{document}

\makemathtitle
    {Regular Languages}
    {A brief introduction to alphabets, strings and languages}
    {Vivek Kalita}
    {Department of Physics, IISER Mohali}
    {20 October 2026}

\tableofcontents

\section{Some basic definitions}

Here I introduce some very basic words.
\begin{definition}[Alphabet, string]
The following are two fundamental building blocks of formal language theory:
\begin{enumerate}
    \item \textcolor{red!60!black}{\textsf{Alphabet}}: An alphabet $\Sigma$ is a finite, non-empty set of symbols, whose elements are called letters.
    \item \textcolor{red!60!black}{\textsf{String}}: A string over an alphabet $\Sigma$ is a finite ordered sequence of elements of $\Sigma$. 
    \begin{enumerate}
        \item The length of a string $w$, denoted by $|w|$, is the number of symbols in the sequence.
        \item The unique string of length zero is called the empty string and is denoted by $\varepsilon$.
    \end{enumerate}
     
\end{enumerate}
\end{definition}
Here are some basic examples: 
\begin{enumerate}
    \item \textcolor{violet!80!black}{\textsf{Binary alphabet.}} Consider the alphabet
    \[
        \Sigma = \{0,1\}.
    \]
    Some strings over this alphabet are
    \[
        0,\quad 1,\quad 001,\quad 1010,\quad 11001.
    \]

    \item \textcolor{violet!80!black}{\textsf{English alphabet.}} Consider the alphabet
    \[
        \Sigma = \{a,b,c,\ldots,z\}.
    \]
    Some strings over this alphabet are
    \[
        \texttt{automata},\quad
        \texttt{theory},\quad
        \texttt{computer}.
    \]

    \item \textcolor{violet!80!black}{\textsf{Decimal alphabet.}} Consider the alphabet
    \[
        \Sigma = \{0,1,2,\ldots,9\}.
    \]
    Some strings over this alphabet are
    \[
        123,\quad 2026,\quad 314159.
    \]
\end{enumerate}

\begin{definition}[Concatenation]
 The \textsf{concatenation} of $u$ and $v$, denoted by $uv$, is the string obtained by appending the symbols of $v$ to the end of $u$. In particular, concatenation satisfies
\[
    |uv|=|u|+|v|,
\]
and the empty string $\varepsilon$ acts as the identity:
\[
    u\varepsilon=\varepsilon u=u.
\]
\end{definition}
The notation $\Sigma^*$ denotes the set of all finite strings over an alphabet $\Sigma$, including the empty string $\varepsilon$.

Let $u,v,w\in\Sigma^*$. The concatenation operation satisfies the following properties:
\begin{enumerate}
    \item \textcolor{violet!70!black}{\textsf{Closure.}} The concatenation of two strings is again a string over $\Sigma$:
    \[
        uv\in\Sigma^*.
    \]

    \item \textcolor{violet!70!black}{\textsf{Associativity.}} Concatenation is associative:
    \[
        (uv)w=u(vw).
    \]

    \item \textcolor{violet!70!black}{\textsf{Identity element.}} The empty string $\varepsilon$ acts as the identity element:
    \[
        u\varepsilon=\varepsilon u=u.
    \]

    \item \textcolor{violet!70!black}{\textsf{Non-commutativity.}} Concatenation is not commutative in general:
    \[
        uv\neq vu.
    \]
\end{enumerate}

It is also important to observe that non-empty strings do not possess inverses under concatenation. This follows immediately from the fact that
\[
    |uv|=|u|+|v|,
\]
which implies that $uv\neq\varepsilon$ whenever $u\neq\varepsilon$. Consequently, the algebraic structure $(\Sigma^*,\cdot)$, where $\cdot$ denotes concatenation, forms a \textsf{monoid}, but not a group.

\subsection{Languages and basic operations}
\begin{definition}[\textsf{Language}]
Let $\Sigma$ be an alphabet. A \textsf{language} $L$ over $\Sigma$ is any subset of $\Sigma^*$, where $\Sigma^*$ denotes the set of all finite strings over $\Sigma$, including the empty string $\varepsilon$. Thus,
\[
    L\subseteq\Sigma^*.
\]
\end{definition}

Some examples of languages are given below:
\begin{enumerate}
    \item \textcolor{violet!70!black}{\textsf{Binary language.}} Consider the alphabet
    \[
        \Sigma=\{0,1\}.
    \]
    The collection of all binary strings ending in $1$ forms a language:
    \[
        L=\{w\in\Sigma^*:w\text{ ends in }1\}.
    \]

    \item \textcolor{violet!70!black}{\textsf{English language.}} The English language can be regarded as a language over an alphabet consisting of the English letters, spaces, and punctuation marks such as full stops, commas, semicolons, and colons. The collection of valid English expressions forms a language over this alphabet.

    \item \textcolor{violet!70!black}{\textsf{Empty language.}} The empty set $\varnothing$ is also a language over any alphabet, since
    \[
        \varnothing\subseteq\Sigma^*.
    \]
\end{enumerate}

\begin{definition}[\textsf{Concatenation of Languages}]
Let $\Sigma$ be an alphabet and let $L_1,L_2\subseteq\Sigma^*$ be two languages. The \textsf{concatenation} of $L_1$ and $L_2$, denoted by $L_1L_2$, is defined as
\[
    L_1L_2=\{uv:u\in L_1,\ v\in L_2\}.
\]
In other words, it is the set of all strings obtained by concatenating a string from $L_1$ with a string from $L_2$.
\end{definition}

Some examples of language concatenation are given below:
\begin{enumerate}
    \item Let
    \[
        L_1=\{0,01\},\qquad L_2=\{1,10\}.
    \]
    Their concatenation is
    \[
        L_1L_2=\{01,010,011,0110\}.
    \]

    \item Let
    \[
        L_1=\{a,ab\},\qquad L_2=\{b,ba\}.
    \]
    Then
    \[
        L_1L_2=\{ab,aba,abb,abba\}.
    \]
\end{enumerate}

The concatenation of languages satisfies the following basic properties. Let $L_1,L_2,L_3\subseteq\Sigma^*$ be languages.
\begin{enumerate}
    \item \textcolor{violet!70!black}{\textsf{Distributivity over union.}} Concatenation distributes over union:
    \[
        L_1(L_2\cup L_3)
        =L_1L_2\cup L_1L_3,
    \]
    \[
        (L_1\cup L_2)L_3
        =L_1L_3\cup L_2L_3.
    \]

    \item \textcolor{violet!70!black}{\textsf{Associativity.}} Concatenation is associative:
    \[
        (L_1L_2)L_3=L_1(L_2L_3).
    \]

    \item \textcolor{violet!70!black}{\textsf{Identity element.}} The singleton language $\{\varepsilon\}$ acts as the identity element:
    \[
        L_1\{\varepsilon\}
        =\{\varepsilon\}L_1=L_1.
    \]
\end{enumerate}

\begin{definition}[\textsf{Kleene Star}]
Let $\Sigma$ be an alphabet and let $L\subseteq\Sigma^*$ be a language. The \textsf{Kleene star} of $L$, denoted by $L^*$, is the set of all strings obtained by concatenating zero or more strings from $L$. Formally,
\[
    L^*=\bigcup_{n=0}^{\infty}L^n,
\]
where
\[
    L^0=\{\varepsilon\},
    \qquad
    L^{n+1}=L^nL.
\]
\end{definition}

Some examples of the Kleene star operator are given below:
\begin{enumerate}
    \item  Consider the language
    \[
        L=\{a\}.
    \]
    Its Kleene star is
    \[
        L^*=\{\varepsilon,a,aa,aaa,\ldots\}.
    \]

    \item Consider the language
    \[
        L=\{0,1\}.
    \]
    Then
    \[
        L^*=\{\varepsilon,0,1,00,01,10,11,\ldots\},
    \]
    which is precisely the set of all finite binary strings.
\end{enumerate}

The Kleene star operator satisfies the following basic properties:
\begin{enumerate}
    \item \textcolor{violet!70!black}{\textsf{Empty string.}} The empty string always belongs to the Kleene star of a language:
    \[
        \varepsilon\in L^*.
    \]

    \item \textcolor{violet!70!black}{\textsf{Containment.}} Every string in $L$ also belongs to $L^*$:
    \[
        L\subseteq L^*.
    \]

    \item \textcolor{violet!70!black}{\textsf{Closure under concatenation.}} The Kleene star of a language is closed under concatenation. That is, if $x,y\in L^*$, then
    \[
        xy\in L^*.
    \]
    \item \textcolor{violet!70!black}{\textsf{Interaction with union.}}
For any two languages $L_1,L_2$,
\[
    (L_1\cup L_2)^*=(L_1^*L_2^*)^*.
\]
\end{enumerate}

\begin{remark}
The Kleene star of the empty language is
\[
    \varnothing^*=\{\varepsilon\}.
\]
Indeed, by definition, $L^0=\{\varepsilon\}$, whereas
$L^n=\varnothing$ for every $n\geq 1$ when $L=\varnothing$.
Thus, the only string in $\varnothing^*$ is the empty string.
\end{remark}
\end{document}